\section{Closing the character-to-window dependency chain}
\label{sec:window-bridge}

\status{WB1--WB8 give one complete implication, from the explicitly
stated arithmetic gap and size data below to the fixed positive
character window and then to RH and Dirichlet GRH. The finite Fourier
identity, ramified-prime terms, subgroup averaging, moving cutoffs,
primitive-character reading, and uniform quantifier order are proved.
The arithmetic gap and the existence of one family satisfying its
premises remain inputs. This section closes the character/analytic
bridge left open in SC2; it does not assert those arithmetic inputs.}

The point of this consolidation is to exhibit the whole chain in one
place. Its leaves are concrete inequalities on one set of data, rather
than a collection of mutually unrelated existence statements. The
source route is \texttt{ChBalance}, \texttt{ChCore}, \texttt{GmAEq},
\texttt{SigALe}, \texttt{CoreGm}, and \texttt{TopWin}.

% BEGIN WB1
\subsection{WB1. The actual characters, including ramification}

Let $A$ be a finite abelian group of order $n$, and write
\[
 G=(\mathbb Z/2\mathbb Z)^2\times A,\quad
 \Phi=\{(0,b,h)\},\quad f=|\Phi|=2n,\quad
 c_0=(1,0,0),\quad c_\infty=(1,1,0).
\]
Use multiplicative notation for characters of the additive group $G$;
the trivial character, written $1$ here, is written $0$ in the Lean
\texttt{AddChar} notation. Put $\chi_d(a,b,h)=(-1)^a$.
An arithmetic datum assigns an inertia subgroup $I_p\le G$ and a
Frobenius class $F_p\in G$ to every prime. Define, for $j\ge1$,
\[
 v_\xi(p^j)=\frac1{|I_p|}\sum_{h\in I_p}\xi(jF_p+h),
 \qquad
 v_\xi(r)=\prod_{p^j\parallel r}v_\xi(p^j)\quad(r\ge1).
 \tag{WB.1}
\]
Thus $|v_\xi(r)|\le1$ and $v_1(r)=1$. Finite-character
orthogonality gives the exact ramified formula
\[
 \frac1f\sum_{\xi\in\widehat G}\overline{\xi(g)}v_\xi(p^j)
 =\frac{2}{|I_p|}\mathbf1_{g-jF_p\in I_p}.
 \tag{WB.2}
\]
Indeed, expand (WB.1) and interchange the two finite sums. The inner
character sum is $|G|=2f$ when $jF_p+h=g$, and zero otherwise.
There is at most one such $h\in I_p$.

Here is why these are the values of the \emph{primitive} Dirichlet
characters on an actual host. Suppose its positive modulus is $N_0$
and its specified reciprocity map
\[
 r:(\mathbb Z/N_0\mathbb Z)^\times\longrightarrow G
\]
is onto. At a prime $p\mid N_0$, use CRT to split the unit group
into its $p$-primary factor and its prime-to-$p$ factor. Let $I_p$
be the image of the first factor and let $F_p$ be the image of a unit
which is $p$ modulo the prime-to-$p$ part. Its value at the
$p$-primary part is immaterial modulo $I_p$. For $p\nmid N_0$,
take $I_p=0$ and $F_p=r(p)$.

Let $\xi^*$ be the primitive inducer of $\xi\circ r$. If
$\xi|_{I_p}$ is nontrivial, character averaging makes (WB.1) zero;
the primitive conductor is then divisible by $p$, and
$\xi^*(p)=0$ as well. If $\xi|_{I_p}=1$, the induced character
factors through the prime-to-$p$ part. Its primitive conductor is
prime to $p$, and its value at $p^j$ is exactly $\xi(jF_p)$.
This proves
\[
 v_\xi(r)=\xi^*(r)\qquad(r\ge1),
 \tag{WB.3}
\]
also at ramified primes, not only on integers prime to $N_0$.
Uniqueness of primitive induction and surjectivity of $r$ imply
\[
 \xi^*=\zeta^*\Longrightarrow\xi=\zeta,
 \qquad \xi^*=1\Longleftrightarrow\xi=1.
 \tag{WB.4}
\]
For the first implication, induce both primitive characters back to
$N_0$, and then use surjectivity. In particular no growing host
introduces repeated copies of one primitive target.

Extend the already centered principal ramp to every $y>0$ by the
same finite-sum formula
\[
 G(y)=\sum_{r\le y}\Lambda(r)(y/r-1)-y\log y+(1+\gamma)y.
\]
For $0<y<1$ the sum is empty. This extension matters in the exact
balance identity, where $a/w$ need not be at least $1$.
Fix numbers $\ell_\xi\in\mathbb C$ for $\xi\ne1$ and set, for
$y>0$,
\[
 \mathcal G_1(y)=G(y),\qquad
 \mathcal G_\xi(y)=\sum_{r\le y}\Lambda(r)v_\xi(r)(y/r-1)
                         +y\ell_\xi\quad(\xi\ne1).
 \tag{WB.5}
\]
On the original host require the actual equality
$\ell_\xi=L'/L(1,\xi^*)$. Equations (WB.3)--(WB.5) then identify
$\mathcal G_\xi$ with the already centered primitive ramp
$G_{\xi^*}$. This is the literal content of the
\texttt{HostAddCharCompatible}/\texttt{prim\_read} join.
The arbitrary value $\ell_1$ is never used.
% END WB1

% BEGIN WB2
\subsection{WB2. Exact balance, with the exceptional prime retained}

Assume that a finite prime set $S$ and a prime $d\notin S$ satisfy
\[
 I_p=0\ (p\notin S\cup\{d\}),\quad
 I_p\subset\Phi,\ F_p\in\Phi\ (p\in S),\quad
 I_d=\{0,c_0\}.
 \tag{WB.6}
\]
These are the substantive fields of \texttt{HostFactsS}. The source
also includes $2,q\in S$ to identify its inert-prime indexing.
For $t\in A$, the map
\[
 g\longmapsto\rho_t(g):=c_\infty+g+(0,0,t)
\]
is a bijection $G\setminus\Phi\to\Phi$. Given $a>0$ and
$w:\Phi\to\mathbb R_{>0}$, choose an integer $N$ at least every
$a/w(\rho)$. Write $(z)_+=\max(z,0)$ and define the actual asked
prime sum
\begin{align*}
 \mathcal A(a,t,w)={}&
 2\sum_{\substack{p\notin S\cup\{d\}\\F_p\notin\Phi}}
    \log p\sum_{\substack{j\ge1\\jF_p\notin\Phi}}
       (a/p^j-w(\rho_t(jF_p)))_+\\
 &+\log d\sum_{j\ge1}
       \sum_{\substack{g\notin\Phi\\g-jF_d\in I_d}}
          (a/d^j-w(\rho_t(g)))_+ .
 \tag{WB.7}
\end{align*}
All sums are finite: terms with $p^j>N$ vanish. This is the source's
\texttt{ASK}; when $d>N$ its last summand vanishes, rather than
requiring a separate assumption on $d$.

Define
\begin{align*}
 \mathcal M(a,w)&=a\left(\log a-1-\gamma+\ell_{\chi_d}
                       -\frac1f\sum_{\rho\in\Phi}\log w(\rho)\right),\\
 \mathcal D(a,t,w)&=\frac1f\sum_{g\notin\Phi}
       \sum_{\xi\in\widehat G}\overline{\xi(g)}
       w(\rho_t(g))\mathcal G_\xi(a/w(\rho_t(g))).
\end{align*}
Then the precise complex identity is
\[
                 \boxed{\ \mathcal A-\mathcal M=\mathcal D.\ }
 \tag{WB.8}
\]
To prove it, apply (WB.2) to each prime power. At an unramified prime
the coefficient is $2$ at $g=jF_p$ and zero elsewhere. At $p\in S$
the coset $jF_p+I_p$ lies in $\Phi$, so it contributes nothing.
At $d$ the coefficient is $1$ on its inertia coset; that coset has
exactly one member outside $\Phi$, because $c_0\notin\Phi$.
Consequently (WB.7) equals
\[
 \frac1f\sum_{g\notin\Phi}\sum_\xi\overline{\xi(g)}
       \sum_{r\ge1}\Lambda(r)v_\xi(r)(a/r-w(\rho_t(g)))_+ .
\]
For $W>0$, the innermost sum is $W\mathcal G_\xi(a/W)-a\ell_\xi$
if $\xi\ne1$. For $\xi=1$ it is
\[
 WG(a/W)+a\log a-a\log W-(1+\gamma)a.
\]
Finally orthogonality on the index-two subgroup $\Phi$ gives
\[
 \sum_{g\notin\Phi}\overline{\xi(g)}=
 \begin{cases} f,&\xi=1,\\-f,&\xi=\chi_d,\\0,&\text{otherwise}.
 \end{cases}
\]
The $-a\ell_{\chi_d}$ correction is therefore multiplied by $-f$:
this proves the \emph{plus} sign in $a\ell_{\chi_d}$ in
$\mathcal M$. The bijection $\rho_t$ gives the logarithmic sum.
This proves (WB.8) without dropping any ramified contribution.
% END WB2

% BEGIN WB3
\subsection{WB3. Uniform regularity of the centered ramps}

The following estimates use only $|v_\xi(r)|\le1$ and elementary
Chebyshev bounds; their constants are independent of the host.
First,
\[
 \Psi(y):=\sum_{r\le y}\Lambda(r)\le3y\quad(y\ge0),\qquad
 \log k-1\le\sum_{r\le k}\frac{\Lambda(r)}r\le\log k+3
 \quad(k\in\mathbb Z_{\ge1}).
 \tag{WB.9}
\]
For completeness, central binomial coefficients give
$\vartheta(y)=\sum_{p\le y}\log p\le y\log4$ by induction on
$\lfloor y\rfloor$. At $2m$ use
$\vartheta(2m)-\vartheta(m)\le\log\binom{2m}{m}\le2m\log2$;
at $2m+1$ use
$\vartheta(2m+1)-\vartheta(m+1)\le
\log\binom{2m+1}{m}\le2m\log2$.
The induction starts at $0,1,2$.
Summing prime powers gives, for $y\ge1$,
\[
 \Psi(y)\le y\log4+2\sqrt y\log y
 \le(\log4+4/e)y<3y;
\]
here $\log t/t\le1/e$ with $t=\sqrt y$. For the second bound use
\[
 \log(k!)=\sum_{r\le k}\Lambda(r)\lfloor k/r\rfloor,
 \qquad k\log k-k\le\log(k!)\le k\log k.
\]
The two inequalities $k/r-1\le\lfloor k/r\rfloor\le k/r$ and
the first part of (WB.9) give the result.

There is an absolute $C_0>0$ such that, for $y,z\ge1$,
\begin{align*}
 |G(y)-G(z)|&\le C_0|y-z|,\\
 |\mathcal G_\xi(y)-\mathcal G_\xi(z)|
 &\le |y-z|\bigl(\log\max(y,z)+C_0+|\ell_\xi|\bigr)
                 \quad(\xi\ne1),\\
 |\mathcal G_\xi(x)|&\le x(\log x+C_0+|\ell_\xi|)
                 \quad(x\ge1,\ \xi\ne1).
 \tag{WB.10}
\end{align*}
Between consecutive integers, differentiate the finite sums. For $G$
the derivative is $\sum_{r\le y}\Lambda(r)/r-\log y+\gamma$,
bounded by (WB.9); replacing $\log\lfloor y\rfloor$ by $\log y$
costs at most $\log2$. For a nontrivial character, bound the
derivative termwise using (WB.9). Each ramp is continuous at the
integers because the newly appearing term is zero there; integrate
the derivative bounds interval by interval. The final estimate uses
$\mathcal G_\xi(1)=\ell_\xi$. One may take $C_0=6+|\gamma|$.

For every character, $a>0$ and $w,w'>0$, we also have the exact
constant needed by the source:
\[
 \left|w\mathcal G_\xi(a/w)-w'\mathcal G_\xi(a/w')\right|
 \le\frac{3a}{\min(w,w')}|w-w'|.
 \tag{WB.11}
\]
For $\xi\ne1$, rewrite the left expression using
\[
 w\mathcal G_\xi(a/w)
 =\sum_{r\ge1}\Lambda(r)v_\xi(r)(a/r-w)_++a\ell_\xi.
\]
The last term cancels; the positive-part function is $1$-Lipschitz,
and (WB.9) bounds the remaining coefficient sum by
$3a/\min(w,w')$. For the trivial character suppose $w\le w'$.
The difference of the two uncentered ramp sums lies between $0$ and
$3a(w'-w)/w$. The centering subtracts
$a\log(w'/w)$, which lies between $0$ and $a(w'-w)/w$.
The difference of these two nonnegative numbers has absolute value
at most $3a(w'-w)/w$. Interchanging $w,w'$ proves (WB.11).
% END WB3

% BEGIN WB4
\subsection{WB4. Subgroup averaging produces the scalar coefficient}

Choose a subgroup $A_0\le A$, sets $T\subset A$ and
$T_2\subset A_0$, and an integer $m\ge1$. Write
\[
 n_T=|T|>0,\quad n_2=|T_2|>0,\quad
 d_2=n_2/n,\quad Q=n_Tmb,\quad
 \mathcal B=\{\xi\in\widehat G:\xi(0,0,A_0)=1\}.
\]
Let $x\ge3$, $b,b',X'>0$ and $\eta,\theta\ge0$. Supply
$a_t>0$ for $t\in T$ and positive weights
$w_{u,j}(\rho)$ for $u\in A_0$, $1\le j\le m$, $\rho\in\Phi$,
with
\begin{gather*}
 |a_t/b-x|\le1,\qquad b'X'\le xb,\\
 b\le w_{u,j}(0,1,s+u)\le b+\eta+1\quad(s\in T_2),\\
 w_{u,j}(\rho)\le2b'\quad
       \bigl(\rho\notin\{(0,1,s+u):s\in T_2\}\bigr),\\
 |\mathcal G_\xi(y)|\le y\theta
       \quad(\xi\in\widehat G,\ y\ge X'/3),\\
 2f(x+1)\theta\le1/x,\qquad
 6f(x+1)(\eta+1)\le b/x.
 \tag{WB.12}
\end{gather*}
The last two inequalities are precisely \texttt{CostWindows}.
The bound for every $y\ge X'/3$ is the supplied
\texttt{ThinScales}, not an estimate inferred from Fourier
orthogonality. Choose a common finite cutoff $N$ for these weights.

The one remaining arithmetic inequality at this point is
\[
 \sum_{t\in T}\sum_{j=1}^m
   \bigl(\mathcal A(a_t,t,w_{u,j})
                     -\Re\mathcal M(a_t,w_{u,j})\bigr)
 \le E\qquad\text{for every }u\in A_0.
 \tag{WB.13}\label{eq:wb-arithmetic-gap}
\]
It must hold for these exact data. It is not supplied by the following
argument.

Put $e_1(t)=(1,0,t)$ and define the ideal thick expression
\[
 \mathcal T_u(t,y)=\frac b f\sum_{s\in T_2}\sum_\xi
             \xi(-e_1(s+u-t))\mathcal G_\xi(y).
\]
The preimage under $\rho_t$ of $(0,1,s+u)$ is
$e_1(s+u-t)$. At such a thick point, (WB.11) bounds the change
from $w\mathcal G_\xi(a_t/w)$ to
$b\mathcal G_\xi(a_t/b)$ by $3(x+1)(\eta+1)$.
At a thin point,
\[
 \frac{a_t}{w}\ge\frac{(x-1)b}{2b'}\ge X'/3
\]
because $x\ge3$ and $b'X'\le xb$. Thus its discarded term has
absolute value at most $a_t\theta\le(x+1)b\theta$.
There are $f$ outer points and $2f$ characters; after division by
$f$, the total error is at most
\[
 2f\{3(x+1)(\eta+1)+(x+1)b\theta\}\le2b/x.
\]
Use (WB.8) and sum over $t,j$ in (WB.13) to obtain
\[
 \sum_{t,j}\Re\mathcal T_u(t,a_t/b)\le E+2Q/x.
 \tag{WB.14}
\]
Now average over $u\in A_0$. If $\xi\notin\mathcal B$ its
subgroup sum vanishes. If $\xi\in\mathcal B$, the factors from
$s$ and $u$ are one, leaving $\xi(e_1(t))$. Consequently
\[
 \frac1{|A_0|}\sum_u\mathcal T_u(t,y)
 =\frac{bd_2}{2}\sum_{\xi\in\mathcal B}
                         \xi(e_1(t))\mathcal G_\xi(y).
 \tag{WB.15}
\]
Combining (WB.14)--(WB.15) proves the complete averaging inequality
\[
 \frac{mbd_2}{2}\sum_{t\in T}
     \Re\sum_{\xi\in\mathcal B}\xi(e_1(t))
                                     \mathcal G_\xi(a_t/b)
 \le E+2Q/x.
 \tag{WB.16}
\]
In particular the factor $d_2/2$ is forced by
$f=2|A|$; there is no extra factor depending on $|A_0|$.
% END WB4

% BEGIN WB5
\subsection{WB5. The biased set gives a fixed positive signal}

Fix $\varepsilon>0$. Let $K\subset\mathcal B$ contain $1$, with $\xi(c_0)=1$ for
$\xi\in K$. Let $0<c\le1$ and define
\[
 \omega_\xi=\begin{cases}1,&\xi(0,0,h)=1\text{ for all }h\in A,\\
                         c,&\text{otherwise}.
              \end{cases}
\]
Assume that $M,K_0,C_E\ge0$ are fixed independently of the growing
host and that
\begin{gather*}
 |\mathcal B|\le M,\qquad
 \left|\sum_{t\in T}\xi(0,0,t)-\omega_\xi n_T\right|
       \le Mx^{-1+\varepsilon}n_T\quad(\xi\in K),\\
 |\ell_\xi|\le K_0\quad(\xi\in K\setminus\{1\}),\qquad
 |\ell_\xi|\le C_E\log x\quad(\xi\in\mathcal B\setminus K).
 \tag{WB.17}
\end{gather*}
For characters trivial on $A$ the first error is exactly zero.
Write
\[
 V(x)=\sum_{\xi\in K}\omega_\xi\Re\mathcal G_\xi(x)
                    -\sum_{\xi\in\mathcal B\setminus K}
                                      |\mathcal G_\xi(x)|.
 \tag{WB.18}
\]
We claim that, for all sufficiently large $x$ with a threshold depending
only on $M,K_0,C_E,\varepsilon$,
\[
 \frac1{n_T}\sum_{t\in T}\Re\sum_{\xi\in\mathcal B}
       \xi(e_1(t))\mathcal G_\xi(a_t/b)
 \ge V(x)-4x^{1/2+\varepsilon}.
 \tag{WB.19}
\]
Here are all the error terms. Put $y_t=a_t/b$. The principal character
costs at most $C_0$ by (WB.10). For a retained nonprincipal character,
$\xi(e_1(t))=\xi(0,0,t)$, and the difference from
$\omega_\xi n_T\mathcal G_\xi(x)$ is
\[
 \sum_t\xi(0,0,t)(\mathcal G_\xi(y_t)-\mathcal G_\xi(x))
 +\left(\sum_t\xi(0,0,t)-\omega_\xi n_T\right)
                                                   \mathcal G_\xi(x).
\]
After division by $n_T$ its norm is at most
\[
 \log(x+1)+C_0+K_0
       +Mx^\varepsilon(\log x+C_0+K_0).
\]
For an omitted character, use absolute values; its lower bound is
$-|\mathcal G_\xi(x)|-\log(x+1)-C_0-C_E\log x$.
Thus the total loss is bounded by $MR(x)$, where $R$ is the maximum
of these three displayed bounds. Since all their coefficients are
fixed,
\[
 \frac{MR(x)}{x^{1/2+\varepsilon}}\longrightarrow0.
\]
This proves (WB.19), including uniformity in $n$, $|A_0|$, the prime
pair, $m$, the individual weights, and the number of all characters
of $G$. Only the bounded quotient family $\mathcal B$ is counted.

The actual biased-set hypothesis supplies the Fourier estimate rather
than assuming it separately. Let $[A:A_0]=o$ and
$c=(3o+1)^{-1}$. If for every coset $r+A_0$
\[
 \left|\#(T\cap(r+A_0))-
       \frac{3+\mathbf1_{r\in A_0}}{3o+1}\,n_T\right|
       \le x^{-1+\varepsilon}n_T,
 \tag{WB.20}
\]
then every nontrivial character of $A/A_0$ has sum over $T$ within
$ox^{-1+\varepsilon}n_T$ of $cn_T$: the constant part vanishes
by orthogonality and the distinguished coset contributes $cn_T$.
The triangle inequality bounds the $o$ errors. The quotient
$G/(0,0,A_0)$ has order $4o$, so
$|\mathcal B|=4o$. Taking $M=4o$ proves (WB.17). This is the
mathematical reason for the source's particular $3:4$ bias.
% END WB5

% BEGIN WB6
\subsection{WB6. Closing the scalar and primitive-character joins}

Combining (WB.16) and (WB.19) gives
\[
 \frac{Qd_2}{2}\bigl(V(x)-4x^{1/2+\varepsilon}\bigr)
                   \le E+2Q/x.
 \tag{WB.21}
\]
Assume now the exact scalar sizes from SC1:
\begin{gather*}
 \frac{n_T}{n}\ge\tfrac12x^{-\varepsilon/8},\qquad
 x^{-\varepsilon}\le d_2\le1,\\
 x^{1/2-\varepsilon}\le m\le x^{1/2+\varepsilon/2},\qquad
 xm\le k\le n^2x^{1/2-\varepsilon},\qquad
 0\le a_E\le2xb,
\end{gather*}
and put
\[
 B=n_Ta_E(1+m^2/k)
       +m\left(2nmbx^\varepsilon+
                \frac{n_2kbx^\varepsilon}{n(nx^{-\varepsilon})}\right).
 \tag{WB.22}
\]
For some fixed $C\ge0$, the remaining gap bound must give
$E\le x^{C\varepsilon}B$. The complete proof of SC1 applied
to (WB.21) gives, once $26\le x^{3\varepsilon/8}$,
\[
                    V(x)\le x^{1/2+(C+3)\varepsilon}.
 \tag{WB.23}
\]
The elementary proof can be read without a dependency search: divide
$B$ by $Q$. Its four terms are at most $2Z,2Z,4Z,2Z$, respectively,
where $Z=x^{1/2+13\varepsilon/8}$. Divide (WB.21) by $Q/2$,
use $d_2\ge x^{-\varepsilon}$, and obtain
$V\le26x^{1/2+(C+21/8)\varepsilon}$. Absorb $26$ as stated.

For the fixed primitive target $\psi$ of order $o$, retain exactly
the primitive inducers of its powers, read via (WB.3). Require this
reading to be onto that fixed family. Its injectivity follows from
(WB.4). The source's cyclic level has \texttt{psK=none}; its
\texttt{EAtH.ps} clause excludes the auxiliary character
$\chi_q(a,b,h)=(-1)^b$ from $K$.
Any character trivial on $A$ and taking value $1$ at $c_0$ is
either $1$ or $\chi_q$. Thus only the trivial character in $K$ has
weight $1$, and all other retained characters have weight $c$.
Consequently the retained term in (WB.18) is exactly
\[
 H_\psi(x)=G(x)+\frac1{3o+1}
       \sum_{\substack{\chi\text{ primitive inducer of }\psi^j\\\chi\ne1}}
                                             \Re G_\chi(x).
 \tag{WB.24}
\]
This equality uses the actual primitive values and the actual linear
terms at $1$. A mere equality of the number of characters would not
suffice. A general source level can retain $\chi_q$ with weight
$1$ through its \texttt{psK} field; no such extra term occurs for
the cyclic levels used here.

If every omitted member of $\mathcal B$ obeys the supplied estimate
$|\mathcal G_\xi(x)|\le x^{1/2+\varepsilon}$, then
\[
 \Sigma(x):=\sum_{\mathcal B\setminus K}|\mathcal G_\xi(x)|
                  \le4o\,x^{1/2+\varepsilon}.
\]
The pointwise conclusion is therefore
\[
 \boxed{\ H_\psi(x)\le x^{1/2+(C+4)\varepsilon}\ }
 \quad\text{when }x^\varepsilon\ge1+4o.
 \tag{WB.25}
\]
In the source $C=e_D+e_L+3\ge0$, giving exactly
$1/2+(e_D+e_L+7)\varepsilon$. All thresholds used to reach
(WB.25) depend on fixed target constants and $\varepsilon$, before
the varying host and its data are considered.
% END WB6

% BEGIN WB7
\subsection{WB7. The complete input contract at the arithmetic boundary}

The preceding implication has the following explicit contract. For
each fixed primitive target and each sufficiently small
$\varepsilon>0$, supply arbitrarily late windows
$[Y,Y^{3/2}]$. For every $x$ in one such window, supply
\emph{one common datum} satisfying all of the following:
\begin{enumerate}
\item An onto reciprocity reading, its actual inertia and Frobenius,
the split/ramified conditions (WB.6), and the actual values
$\ell_\xi=L'/L(1,\xi^*)$.
\item A subgroup of fixed index $o$, the fixed primitive-power
retained family, its auxiliary-character exclusion, the biased
counts (WB.20), and fixed bounds $K_0,C_E$ in (WB.17).
\item The same $T,T_2,a_t,w_{u,j},m,b,b',X',\eta,\theta,N$
satisfying (WB.12), together with the omitted-ramp estimates in WB6.
\item The scalar size inequalities of WB6 and the actual arithmetic
gap (\ref{eq:wb-arithmetic-gap}) with
$E\le x^{C\varepsilon}B$, for a fixed $C\ge0$.
\end{enumerate}
Here $C$ is chosen before $\varepsilon$. Bounds such as $K_0,C_E$
may be chosen for the fixed target (or even for a fixed
$\varepsilon$), but must be uniform as the windows move to infinity.
The common datum may depend on $x$; the target family and $C$ may not.

Every implication from this contract to (WB.25) is proved in
WB1--WB6. There are no remaining character-orthogonality,
ramified-value, moving-cutoff, or scalar-comparison inputs in this
part of the chain. The unproved task is to construct this contract
from the arithmetic geometry on the same maps and fields.

For precise navigation, the source packages items 1--3 in
\texttt{HostAx}, \texttt{EAtH}, the reading part of
\texttt{CoreAtH}, \texttt{SizesA}, \texttt{ThinScales}, and
\texttt{CostWindows}. Item 4 is the conclusion of \texttt{ch\_top}
and the remaining fields of \texttt{CoreAtH}. Its gap is obtained
there from the eleven-field \texttt{SizeData2} package, uniform
error windows, and the two bounds on $B_0,N_c$. Those are not proved
merely by unpacking \texttt{CoreAtH}.

The algebraic part of that last source join is short and worth
recording. \texttt{IdentU} supplies the exact bijections from the
original top indices to $T$ and from the place indices to $\Phi$,
the equalities of $a,w,\gamma$, the equality
$\mathrm{LpL}=\Re\ell_{\chi_d}$, and the equality of the asked
sums. Substitution into the definition of the ledger main term gives
\[
 \mathrm{Main}(u)=\sum_{t,j}\Re\mathcal M(a_t,w_{\mathrm{tr}(u),j}),
 \qquad
 \mathrm{ASK}(u)-\mathrm{Main}(u)
    =\operatorname{gap}(\mathrm{tr}(u)).
 \tag{WB.26}
\]
The source's coverage clause $A_0\subset\mathrm{range}(\mathrm{tr})$
then transfers its bound to every $u\in A_0$ needed in (WB.13).
The $B_0,N_c$ inequalities replace its error bracket by (WB.22),
termwise and with nonnegative coefficients. Finally
$e_{\rm Div}\le e_D$, $e_{\rm Loss}\le e_L$ bound its error
exponent by $C=e_D+e_L+3$. These are exact substitutions on one
supplied ledger. The independent divisibility, metric, dimension,
theta-norm, and spectral estimates that establish that ledger are
still the arithmetic frontier.
% END WB7

% BEGIN WB8
\subsection{WB8. A closed terminal theorem with the quantifiers visible}

\begin{theorem}[Arithmetic contract implies the fixed target's zero-free half-plane]
Fix a primitive Dirichlet character $\psi$. Suppose there exist
$C\ge0$ and $\varepsilon_0>0$ such that for every
$0<\varepsilon<\varepsilon_0$, the contract in WB7 can be supplied
on arbitrarily late windows $[Y,Y^{3/2}]$. Then every primitive
inducer of a power of $\psi$ has no zero in $\Re s>1/2$,
away from the principal pole. If the contract is available for every
primitive $\psi$, RH and GRH for all Dirichlet $L$-functions follow.
\end{theorem}

\begin{proof}
Choose $0<\delta<1/2$ and then
\[
 0<\varepsilon<\min\{\varepsilon_0,\delta/(C+4)\}.
\]
All the pointwise thresholds in WB3--WB6 are now fixed. Choose one
of the contract's windows whose left endpoint exceeds them and any
prescribed $Y_0$. Equation (WB.25) holds at \emph{every} real $x$
in this window and gives $H_\psi(x)\le x^{1/2+\delta}$.
The positive-family one-sided criterion AN3 now rules out any
constituent zero with real part greater than $1/2+\delta$.
Given a proposed zero with real part $\beta>1/2$, choose
$0<\delta<\min\{1/2,\beta-1/2\}$; this is a contradiction. Positivity of
all coefficients in (WB.24), including the residue at a possible
real zero, is the point that allows a one-sided upper bound.

For an arbitrary character induced by primitive $\chi^*$,
\[
 L(s,\chi)=L(s,\chi^*)
       \prod_{\substack{p\mid N\\p\nmid\mathrm{cond}(\chi^*)}}
                       (1-\chi^*(p)p^{-s}).
\]
Every additional factor is nonzero for $\Re s>0$. The primitive
functional equation maps a zero left of $1/2$ in the open critical
strip to a zero right of $1/2$ of the conjugate character; its gamma,
conductor, and root-number factors are nonzero there. This proves
the all-character strip statement. The ordinary target $\psi=1$
has $o=1$ and $H_\psi=G$, giving RH by AN2, including its standard
functional-equation treatment outside the open strip. At $s=1$,
nonprincipal characters use their classical nonvanishing theorem;
the principal pole is excluded as in the exported statement.
\end{proof}

The retained analytic leaves are exactly those stated in AN2--AN3
and GH6: ordinary analytic continuation, the functional equations,
Euler products and classical nonvanishing on $\Re s=1$, together
with elementary complex analysis. This theorem adds no cohomology
or representation hypothesis beyond WB7's explicitly open
arithmetic contract.

The source's \texttt{tiedWin\_of\_ex} supplies the matching
quantifier rule: first choose the common bounds demanded by all
links, and only then use a family-existence theorem valid for
\emph{every} such bound. The dependency audit below spells this out.
On the literal \texttt{ShapeWin}\,$\to$\,\texttt{topDesWin}
route the chosen values are $e_D=2,e_L=0$, so its conditional
terminal exponent is $1/2+9\varepsilon$. This is an identification
of the constants in that route, not a new construction of its
required simultaneous family.
% END WB8
